
El experimento consiste en evaluar el rendimiento de los cuatro algoritmos implementados bajo diferentes escenarios. Se generaron 27 casos de prueba divididos en tres categorías:
\begin{itemize}
    \item \textbf{Casos Pequeños ($n \in \{3, 5, 8\}$):} Diseñados para validar que la Fuerza Bruta y la Programación Dinámica coinciden en el resultado óptimo.
    \item \textbf{Casos Medianos ($n \in \{20, 40, 80\}$):} Orientados a medir la brecha de calidad entre las heurísticas Greedy y la solución óptima de DP.
    \item \textbf{Casos Grandes ($n \in \{100, 150, 200\}$):} Utilizados para analizar el escalamiento temporal y de memoria de los algoritmos eficientes.
\end{itemize}

Las pruebas se ejecutaron en una máquina con sistema operativo Linux, midiendo el tiempo de ejecución en milisegundos mediante la librería \texttt{chrono} y el uso de memoria máximo (\texttt{ru\_maxrss}) mediante \texttt{getrusage}.
